\chapter{Elephantiasis}

\section*{Definition and Overview}
Elephantiasis is the massive swelling of a limb or other body part, most commonly the legs and genitals, caused by blockage of the lymphatic system. The most common cause worldwide is lymphatic filariasis, a parasitic infection spread by mosquitoes, in which thread-like worms block the lymph vessels. The term can also describe severe swelling from other causes of lymphatic damage, including surgery and radiotherapy for cancer. The condition causes disfigurement, disability, and stigma, but it is preventable, and modern management greatly improves symptoms and quality of life.

\section*{Epidemiology}
Lymphatic filariasis affects tens of millions of people worldwide, mostly in tropical and subtropical regions of Africa, Asia, the Pacific, and the Americas, making it one of the leading causes of permanent disability in developing countries. In Australia, the disease is rare, although it was historically present in parts of northern Australia. Elephantiasis from other causes, such as cancer treatment or surgery, occurs in Australia and requires similar management.

\section*{Aetiology and Pathophysiology}
Elephantiasis results from damage to the lymphatic system, which normally drains fluid from tissues.
\begin{itemize}
    \item \textbf{Lymphatic filariasis:} caused by parasitic worms, mainly Wuchereria bancrofti, transmitted by mosquito bites; the adult worms live in the lymph vessels and nodes, causing inflammation and blockage
    \item \textbf{Secondary lymphoedema:} lymphatic damage from surgery, particularly for cancer, radiotherapy, injury, or recurrent infection
    \item \textbf{Primary lymphoedema:} congenital or inherited abnormalities of the lymphatics
    \item \textbf{Mechanism:} blocked lymphatics fail to drain fluid from tissues, which accumulates and causes swelling; over time, the tissue becomes thick, hard, and fibrous
    \item \textbf{Recurrent infection:} damaged skin and lymphatics are prone to bacterial infection, which worsens the swelling
    \item \textbf{The name:} the limb can become hugely enlarged and skin thickened like an elephant's, giving the condition its name
\end{itemize}

\section*{Clinical Features}
Elephantiasis develops slowly over years.
\begin{itemize}
    \item Progressive swelling of the leg, arm, scrotum, or breast
    \item A heavy, tight, or aching feeling in the affected part
    \item Thickened, hard, and warty skin
    \item Recurrent episodes of fever and infection of the skin and lymphatics
    \item Difficulty walking and carrying out daily activities when the legs are affected
    \item In men, swelling of the scrotum; in women, swelling of the vulva or breast
    \item Secondary changes: skin folds, ulcers, and foul odour from infection
\end{itemize}

\section*{Differential Diagnosis}
Other causes of limb and genital swelling must be distinguished.
\begin{itemize}
    \item Chronic venous insufficiency and varicose veins
    \item Heart, kidney, and liver disease causing generalised swelling
    \item Deep vein thrombosis
    \item Obesity-related limb swelling
    \item Tumours compressing the lymphatic or venous system
    \item Other causes of lymphoedema, including surgery and radiotherapy
\end{itemize}

\section*{When to Seek Medical Care}
People with persistent or progressive swelling of a limb or genital area should see a doctor for diagnosis. Anyone living in or travelling from areas where filariasis occurs should seek advice if swelling develops, as should people with lymphoedema after cancer treatment.

\textbf{Call 000 (or local emergency services) for:}
\begin{itemize}
    \item Sudden swelling of one leg with pain, which may indicate a blood clot
    \item Sudden severe pain, redness, and fever, which may indicate a serious infection
    \item Difficulty breathing or signs of fluid overload
    \item Fever with a rapidly spreading skin infection
\end{itemize}

\section*{Diagnosis and Investigations}
Diagnosis identifies the cause of the lymphatic blockage.
\begin{itemize}
    \item \textbf{History:} travel history, cancer treatment, and the pattern of swelling
    \item \textbf{Examination:} the extent, distribution, and nature of the swelling
    \item \textbf{Blood tests:} blood films tested at night can detect filarial worms; antibody and antigen tests are also used
    \item \textbf{Ultrasound:} can visualise the worms and assess lymphatic damage
    \item \textbf{Lymphoscintigraphy:} an imaging test of lymphatic drainage
    \item \textbf{Exclusion of other causes:} tests for heart, kidney, and venous disease
\end{itemize}

\section*{Management}
Management controls symptoms, prevents progression, and reduces complications.
\begin{itemize}
    \item \textbf{Antiparasitic medicines:} for filariasis, medicines such as diethylcarbamazine and ivermectin kill the microfilariae and reduce transmission
    \item \textbf{Compression therapy:} bandages, stockings, and pneumatic pumps reduce swelling
    \item \textbf{Exercise and elevation:} specific exercises and limb elevation drain the lymphatics
    \item \textbf{Skin care:} meticulous hygiene and moisturising to prevent infection
    \item \textbf{Management of infection:} prompt antibiotics for episodes of lymphangitis and cellulitis
    \item \textbf{Surgery:} for advanced disease, procedures can remove excess tissue or bypass blocked lymphatics
    \item \textbf{Psychosocial support:} addressing stigma, disability, and quality of life
\end{itemize}

\section*{Complications}
\begin{itemize}
    \item Recurrent bacterial infections and cellulitis
    \item Progressive disability and difficulty walking
    \item Disfigurement and social stigma
    \item Psychological distress and depression
    \item Ulceration and skin breakdown
    \item Rarely, development of a lymphatic cancer
\end{itemize}

\section*{Prognosis and Outcomes}
With modern management, the outlook for people with elephantiasis has improved greatly. Mass drug administration programmes have reduced filariasis transmission dramatically worldwide, and good hygiene, compression, and exercise control the swelling and prevent infection. The condition is chronic, but most people maintain function and independence with consistent care.

\section*{Prevention and Self-Care}
\begin{itemize}
    \item Prevent mosquito bites in affected regions: use repellent, nets, and long clothing
    \item Participate in mass drug administration programmes where offered
    \item Care for the skin daily with washing, drying, and moisturising
    \item Elevate the affected limb and perform prescribed exercises
    \item Wear compression garments as advised
    \item Seek prompt treatment of any skin infection
    \item Maintain a healthy weight, which reduces the burden on the lymphatics
\end{itemize}

\section*{Sources and Further Reading}
The following reputable sources provide further information for healthcare students and patients seeking reliable, current advice about elephantiasis and lymphoedema.
\begin{itemize}
    \item World Health Organization. \textit{Lymphatic filariasis fact sheet.} https://www.who.int/
    \item Healthdirect Australia. \textit{Lymphoedema.} https://www.healthdirect.gov.au/lymphoedema
    \item Australian Lymphoedema Association. \textit{Lymphoedema information and support.} https://www.lymphoedema.org.au/
    \item Centers for Disease Control and Prevention. \textit{Lymphatic filariasis.} https://www.cdc.gov/parasites/lymphaticfilariasis/
    \item MedlinePlus. \textit{Elephantiasis.} https://medlineplus.gov/ency/article/001088.htm
    \item MSD Manual. \textit{Lymphatic filariasis.} https://www.msdmanuals.com/
\end{itemize}
